% GENERATED FILE. Edit canonical lessons, then run npm run generate:books.
% canonical-source-hash: fnv1a64:2da40087e8223941
% canonical-lessons: GE-C117-geldboerse, GE-C117-ring, GE-C117-kette, GE-C117-geschenk, GE-C117-foto

\chapter{Purse, Ring, Chain, Present, Photo}
\label{ch:ge-purse-ring-chain-present-photo}
\hlchaptermodality{117}

\begin{chapteropening}
\textbf{By the end of this chapter:} \emph{I can say a purse, a ring, a chain, a present and a photo.}

Complete the last lesson of chapter 117: I can say a purse, a ring, a chain, a present and a photo.
\end{chapteropening}

\section[die Geldbörse]{die Geldbörse — a purse, a wallet}

\label{lesson:GE-C117-geldboerse}



\begin{quote}
Before the new one: say the German for glasses, then the German for an umbrella.

Say \emph{der Abend}. Then: which of \emph{der, die, das} goes with \emph{Geldbörse}? (\emph{die}.)
\end{quote}

\subsection*{You'll want to know: die Geldbörse}
\textbf{die Geldbörse} — \textquotedblleft{}a purse, a wallet\textquotedblright{}.

\begin{cousinweb}
\textbf{Geldbörse} is \emph{Geld}, \textquotedblleft{}money\textquotedblright{}, + \emph{Börse}, \textquotedblleft{}purse\textquotedblright{}, which came through Dutch and Old French from Late Latin \emph{bursa}, \textquotedblleft{}leather bag\textquotedblright{} (EtymDB 2.1). English \emph{purse} shares it.
\end{cousinweb}

\begin{grammarlens}[title={What you've built}]
One more word for this chapter.
\end{grammarlens}

\subsection*{Your turn}
\begin{itemize}
\raggedright
  \item \emph{Say it:} \emph{die Geldbörse}
  \item \emph{Say it:} \emph{die Geldbörse}, once more
  \item \emph{Say it:} say \emph{die Geldbörse}
  \item \emph{Recall:} say the German for glasses, then the German for an umbrella, then say \emph{die Geldbörse} again
\end{itemize}

\begin{tcolorbox}[breakable,colback=teal!4,colframe=teal!35!black,title={Before you move on}]
What does die Geldbörse mean? (\textquotedblleft{}A purse, a wallet\textquotedblright{}.) Say it once more.
\end{tcolorbox}
\section[der Ring]{der Ring — a ring}

\label{lesson:GE-C117-ring}



\begin{quote}
Before the new one: say the German for an umbrella, then the German for a purse.
\end{quote}

\subsection*{You'll want to know: der Ring}
\textbf{der Ring} — \textquotedblleft{}a ring\textquotedblright{}.

\begin{cousinweb}
\textbf{Ring} $\leftarrow$ Old High German \emph{ring} — inherited. EtymDB 2.1 traces English \emph{ring}, Old English \emph{hring}, to Proto-Germanic \emph{*hringaz}, \textquotedblleft{}something bent\textquotedblright{}: the same word.
\end{cousinweb}

\begin{grammarlens}[title={What you've built}]
One more word for this chapter.
\end{grammarlens}

\subsection*{Your turn}
\begin{itemize}
\raggedright
  \item \emph{Say it:} \emph{der Ring}
  \item \emph{Say it:} \emph{der Ring}, once more
  \item \emph{Say it:} say \emph{der Ring}
  \item \emph{Recall:} say the German for an umbrella, then the German for a purse, then say \emph{der Ring} again
\end{itemize}

\begin{tcolorbox}[breakable,colback=teal!4,colframe=teal!35!black,title={Before you move on}]
What does der Ring mean? (\textquotedblleft{}A ring\textquotedblright{}.) Say it once more.
\end{tcolorbox}
\section[die Kette]{die Kette — a chain; a necklace}

\label{lesson:GE-C117-kette}



\begin{quote}
Before the new one: say the German for a purse, then the German for a ring.
\end{quote}

\subsection*{You'll want to know: die Kette}
\textbf{die Kette} — \textquotedblleft{}a chain; a necklace\textquotedblright{}.

\begin{cousinweb}
\textbf{Kette} $\leftarrow$ Middle High German \emph{keten} (EtymDB 2.1), from Latin \emph{catēna}, \textquotedblleft{}chain\textquotedblright{} — borrowed. English \emph{chain} comes from the same Latin word through French.
\end{cousinweb}

\begin{grammarlens}[title={What you've built}]
One more word for this chapter.
\end{grammarlens}

\subsection*{Your turn}
\begin{itemize}
\raggedright
  \item \emph{Say it:} \emph{die Kette}
  \item \emph{Say it:} \emph{die Kette}, once more
  \item \emph{Say it:} say \emph{die Kette}
  \item \emph{Recall:} say the German for a purse, then the German for a ring, then say \emph{die Kette} again
\end{itemize}

\begin{tcolorbox}[breakable,colback=teal!4,colframe=teal!35!black,title={Before you move on}]
What does die Kette mean? (\textquotedblleft{}A chain; a necklace\textquotedblright{}.) Say it once more.
\end{tcolorbox}
\section[das Geschenk]{das Geschenk — a present}

\label{lesson:GE-C117-geschenk}



\begin{quote}
Before the new one: say the German for a ring, then the German for a chain; a necklace.
\end{quote}

\subsection*{You'll want to know: das Geschenk}
\textbf{das Geschenk} — \textquotedblleft{}a present\textquotedblright{}.

\begin{cousinweb}
\textbf{Geschenk} is \emph{ge-} + the stem of the verb \textquotedblleft{}to give as a present\textquotedblright{}, which first meant \textquotedblleft{}to pour out a drink for someone\textquotedblright{}. A present began as a drink poured for a guest. \emph{Schenken} $\leftarrow$ Old High German \emph{scenken}, \textquotedblleft{}to pour, to give drink\textquotedblright{} (EtymDB 2.1).
\end{cousinweb}

\begin{grammarlens}[title={What you've built}]
One more word for this chapter.
\end{grammarlens}

\subsection*{Your turn}
\begin{itemize}
\raggedright
  \item \emph{Say it:} \emph{das Geschenk}
  \item \emph{Say it:} \emph{das Geschenk}, once more
  \item \emph{Say it:} say \emph{das Geschenk}
  \item \emph{Recall:} say the German for a ring, then the German for a chain; a necklace, then say \emph{das Geschenk} again
\end{itemize}

\begin{tcolorbox}[breakable,colback=teal!4,colframe=teal!35!black,title={Before you move on}]
What does das Geschenk mean? (\textquotedblleft{}A present\textquotedblright{}.) Say it once more.
\end{tcolorbox}
\section[das Foto]{das Foto — a photo}

\label{lesson:GE-C117-foto}



\begin{quote}
Before the new one: say the German for a chain; a necklace, then the German for a present.
\end{quote}

\subsection*{You'll want to know: das Foto}
\textbf{das Foto} — \textquotedblleft{}a photo\textquotedblright{}.

\begin{cousinweb}
\textbf{Foto} is short for the word for \textquotedblleft{}photography\textquotedblright{}, a modern coinage from Greek \emph{phōs}, \textquotedblleft{}light\textquotedblright{}, and \emph{gráphein}, \textquotedblleft{}to write\textquotedblright{}: writing with light. EtymDB 2.1 has no source for \emph{Foto}.
\end{cousinweb}

\begin{grammarlens}[title={What you've built}]
One more word for this chapter.
\end{grammarlens}

\subsection*{Your turn}
\begin{itemize}
\raggedright
  \item \emph{Say it:} \emph{das Foto}
  \item \emph{Say it:} \emph{das Foto}, once more
  \item \emph{Say it:} say \emph{das Foto}
  \item \emph{Recall:} say the German for a chain; a necklace, then the German for a present, then say \emph{das Foto} again
\end{itemize}

\begin{tcolorbox}[breakable,colback=teal!4,colframe=teal!35!black,title={Before you move on}]
What does das Foto mean? (\textquotedblleft{}A photo\textquotedblright{}.) Say it once more.
\end{tcolorbox}
